\documentclass[10pt,a4paper]{amsart}
\usepackage[margin=19mm]{geometry}
\usepackage{amsmath,amssymb,mathtools}
\usepackage[scr=rsfs,cal=euler]{mathalfa}
\usepackage{xfrac}
\usepackage[unicode,hidelinks,bookmarks=false]{hyperref}
\newcommand{\ie}{i.e.\ }
\newcommand{\cf}{cf.\ }
\newcommand{\resp}{respectively\ }
\newcommand{\Subsection}{\subsection}
\providecommand{\nobreakdash}{\nobreak\hbox{-}}
\setlength{\emergencystretch}{2em}
\multlinegap0pt
\usepackage{bm}
\usepackage{bbm}
\usepackage{dsfont}
\usepackage{xspace}
\usepackage{cancel}
\usepackage{mathtools}
\usepackage{amsthm}
\makeatletter
\@ifundefined{remark}{\newtheorem*{remark}{Remark}}{}
\@ifundefined{theorem}{\newtheorem{theorem}{Theorem}}{}
\@ifundefined{problem}{\newtheorem{problem}[theorem]{Problem}}{}
\makeatother
\DeclareMathOperator{\dist}{dist}
\newcommand{\cF}{\mathcal{F}}
\newcommand{\cU}{\mathcal{U}}
\renewcommand{\tilde}{\widetilde}
\title[A functional inequality related to Domar's uniform boundedne]{A functional inequality related to Domar's uniform boundedness theorem}
\author{Thomas Ransford}
\date{Source: arXiv:2607.08718v2 (CC BY 4.0)}
\begin{document}
\maketitle

\noindent\emph{Source and licence.} A functional inequality related to Domar's uniform boundedness theorem. Version arXiv:2607.08718v2 (CC BY 4.0), distributed under the Creative Commons Attribution 4.0 International licence, as recorded in the arXiv licence metadata for this exact version and shown on the abstract page https://arxiv.org/abs/2607.08718v2. Full rights notice: licenses/arxiv-2607.08718-CC-BY-4.0.txt.\par\smallskip

\noindent\textit{Editorial scope.} Counted unit: the Remark whose source label is \texttt{None} in the article (source0.tex lines 644--647, source label \texttt{None}), delivered together with the article's own complete original proof of it (source0.tex lines 650--773, one \texttt{proof} environment of 124 lines). No mathematics is written, repaired or re-derived: statement and proof are the article's own text line for line. Required context retained uncounted: Pb:fnineq (lines 89--100); T:main (lines 150--174); E:nc (lines 153--155); E:Gdef (lines 171--173); E:phiint (lines 627--629); E:Ubound (lines 631--633); E:Domarineq (lines 637--641). Every definition, display and label of those lines is retained unchanged. Omitted from this package and not redistributed: 1--88, 101--149, 156--170, 174--626, 630--630, 634--636, 642--643, 648--649, 774--946. Nothing inside a retained excerpt is omitted or altered: every retained body is delivered exactly as the article's lines are written, so no omission receipt of any kind accompanies this package. Citations: the retained excerpts carry no citation command at all, so no bibliography entry of the article is transcribed and no reference is redistributed. Reference adaptation: none was needed and none was made; every reference inside the delivered unit resolves inside it, so no unresolved reference or citation is produced. Printed slips kept literal: no printed word, symbol, hypothesis or number of the counted statement or of its proof was altered. Layout and class: the article's own class is \texttt{amsart} with packages including amssymb, amsmath, amsthm, newlfont, enumerate; the wrapper is the lane's pinned \texttt{amsart} editorial wrapper at 10pt on A4, which restates the theorem-like environments the retained text needs and adds the article's own macro definitions that the retained text needs (\texttt{\textbackslash cF}, \texttt{\textbackslash cU}, \texttt{\textbackslash dist}, \texttt{\textbackslash tilde}), with extra wrapper packages (bm, bbm, dsfont, xspace, cancel, mathtools). The delivered numbering is the unit's own. Assets: no retained span contains a figure environment, a graphics-inclusion command, a PDF-page inclusion, a table or a TikZ picture; no image file is copied into this package, the package declares no asset dependency and no image file is redistributed with it. Licence: version arXiv:2607.08718v2 of this article is distributed under the Creative Commons Attribution 4.0 International licence, as recorded in the arXiv licence metadata for this exact version and shown on the abstract page https://arxiv.org/abs/2607.08718v2. A full rights notice is delivered at licenses/arxiv-2607.08718-CC-BY-4.0.txt. Characters: every retained excerpt is pure ASCII, so the delivered engine, XeLaTeX with its default Latin Modern text font, reports no missing character.
\begin{problem}\label{Pb:fnineq}
Given a decreasing function $g:(0,\infty)\to[0,\infty)$  
and a constant $\alpha\in(0,1)$, let $\cF(g,\alpha)$ be the family of all 
decreasing functions
 $f:(0,\infty)\to[0,\infty)$
satisfying
\begin{equation}\label{E:fnineq}
f(r+s)\le g(r)+\alpha f(s) \quad(r,s>0).
\end{equation}
Is $\sup_{f\in\cF(g,\alpha)}f(x)<\infty$ for all $x\in(0,\infty)$? 
If so, then find explicit bounds.
\end{problem}

\begin{theorem}\label{T:main}
Let $g,\alpha$ and $\cF(g,\alpha)$ be as specified in Problem~\ref{Pb:fnineq}.
If
\begin{equation}\label{E:nc}
\int_0^1 \log^+g(x)\,dx<\infty,
\end{equation}
then
\begin{equation}\label{E:supfinite}
\sup_{f\in\cF(g,\alpha)}f(x)<\infty \quad(x\in(0,\infty)).
\end{equation}
For $x>0$, define 
\[
F(x):=\sup_{f\in\cF(g,\alpha)}f(x),
\]
and extend $F$ to $0$ by defining  $F(0):=\lim_{x\to0^+}F(x)$ (possibly $\infty$).
Then $F$ satisfies
\begin{equation}\label{E:mainineq}
F\Bigl(\frac{G(x)}{\log(1/\beta)}\Bigr)
\le \frac{g(x-)}{\beta-\alpha}\quad(x\in I,\,\beta\in(\alpha,1)),
\end{equation}
where $I:=\{x\in(0,\infty):g(x)>0\}$ and
\begin{equation}\label{E:Gdef}
G(x):=\int_0^x \log\frac{g(t)}{g(x)}\,dt \quad(x\in I).
 \end{equation}
\end{theorem}

\begin{equation}\label{E:phiint}
 \int_0^1 \log^+\phi^*(t^p)\,dt<\infty.
\end{equation}

\begin{equation}\label{E:Ubound}
\sup_{u\in\cU(Z,\phi)}u(z)\le \psi\bigl(\dist(z,\partial Z)\bigr) \quad(z\in Z),
\end{equation}

\begin{equation}\label{E:Domarineq}
\psi\Bigl(\frac{1}{\alpha^{1/p}\log(1/\beta)}\Bigl(\frac{V_q}{V_{p+q}}\Bigr)^{1/p}
\int_0^t \log\Bigl(\frac{\phi^*(s^p)}{\phi^*(t^p)}\Bigr)\,ds\Bigr)
\le \frac{\phi^*(t^p-)}{\beta-\alpha}.
\end{equation}

\begin{remark}
The quantities $\alpha,\beta$ are parameters subject only to the constraint that $0<\alpha<\beta<1$.
They may be chosen to optimize the bound on $\psi$.
\end{remark}

\begin{proof}
Let $u\in\cU(Z,\phi)$. We shall estimate $u$.
Replacing $u$ by $\max\{u,0\}$,
we may as well suppose  that $u\ge0$.
For convenience, we also set $n:=p+q$.

Let $R$ be the radius of the largest open ball contained
in $Z$. 
For $t\in(0,R]$, set
\[
\tilde{u}(t):=\sup\{u(z):z\in Z,~\dist(z,\partial Z)\ge t\}.
\]
This is finite, since it is the supremum of an upper
semicontinuous function over a compact set.
Extend $\tilde{u}$ to the whole of $(0,\infty)$
by setting $\tilde{u}\equiv 0$ on $(R,\infty)$.
Clearly $\tilde{u}:(0,\infty)\to[0,\infty)$ is a decreasing function.

Let $t\in(0,R]$ and let $z\in Z$ be a point with
$\dist(z,\partial Z)\ge t$. Let $r\in(0,t)$, and let $B$
be the closed ball with centre $z$ and radius $r$.
Then $B\subset Z$, and by the triangle inequality
we have the simple estimate $u\le \tilde{u}(t-r)$ on~$B$.
By the sub-mean
inequality for subharmonic functions, for each $\lambda>0$ we have
\begin{align*}
u(z)
&\le \frac{1}{m_n(B)}\int_B u\,dm_n\\
&\le \frac{1}{m_n(B)}\int_{B\cap(\{\phi\le \lambda\}\times Y)}u\,dm_n
+ \frac{1}{m_n(B)}\int_{B\cap(\{\phi> \lambda\}\times Y)}u\,dm_n\\
&\le \frac{1}{m_n(B)}\int_{B\cap (\{\phi\le \lambda\}\times Y)}\lambda \,dm_n
+ \frac{1}{m_n(B)}\int_{B\cap (\{\phi> \lambda\}\times Y)}\tilde{u}(t-r)\,dm_n\\
&\le \lambda +\frac{m_n(B\cap (\{\phi>\lambda\}\times Y)}{m_n(B)}\tilde{u}(t-r).
\end{align*}
Now $m_n(B)=V_nr^n$. Also, we have
\[
m_n\bigl(B\cap (\{\phi>\lambda\}\times Y)\bigr)
\le m_p\{\phi>\lambda\}(V_qr^q)
=m_1\{\phi^*>\lambda\}(V_qr^q).
\]
Hence
\[
u(z)\le \lambda+\frac{V_q r^q}{V_nr^n}m_1\{\phi^*>\lambda\}\tilde{u}(t-r)
= \lambda+\frac{V_q}{V_nr^p}m_1\{\phi^*>\lambda\}\tilde{u}(t-r).
\]
We are free to choose $\lambda$ how we please.
Let us take $\lambda:=\phi^*(\alpha V_nr^p/V_q)$, where $\alpha\in(0,1)$.
Then we have
$m_1\{\phi^*>\lambda\}\le \alpha V_nr^p/V_q$, and so
\[
u(z)\le \phi^*(\alpha V_nr^p/V_q)+\alpha \tilde{u}(t-r).
\]
Taking the supremum over all $z\in Z$ with
$\dist(z,\partial Z)\ge t$, we obtain the relation
\[
\tilde{u}(t)\le \phi^*(\alpha V_nr^p/V_q)+\alpha \tilde{u}(t-r).
\]
This is valid for all $t\in(0,R]$. 
It also holds trivially if $t\in(R,\infty)$,
since $\tilde{u}$ is zero on this set.
Thus, finally, we deduce that
\begin{equation}\label{E:relation}
\tilde{u}(r+s)\le \phi^*(\alpha V_nr^p/V_q)+\alpha \tilde{u}(s)
\quad(r,s>0,~\alpha\in(0,1)).
\end{equation}


We now apply Theorem~\ref{T:main} with
\[
g(r):=\phi^*(\alpha V_n r^p/V_q).
\]
The condition \eqref{E:nc} for $g$ 
is equivalent to the condition \eqref{E:phiint} for $\phi^*$.
By \eqref{E:relation}, each function $\tilde{u}$ with $u\in\cU(Z,\phi)$ belongs to $\cF(g,\alpha)$. Hence, if 
$F$ denotes the function defined in
Theorem~\ref{T:main}, then
$\tilde{u}(s)\le F(s)$ for all $s>0$.
Consequently, if we define
\[
\psi(s):=\sup_{u\in\mathcal U(Z,\phi)}\widetilde u(s) \qquad(s>0),
\]
and extend $\psi$ to $0$ by setting $\psi(0):=\lim_{s\to0^+}\psi(s)$,
then
\[
\psi(s)\le F(s)\qquad(s\ge0).
\]
In particular, Theorem~\ref{T:main} gives
\[
\psi\Bigl(\frac{G(r)}{\log(1/\beta)}\Bigr)
\le \frac{g(r-)}{\beta-\alpha}
\quad(g(r)>0,\ 0<\alpha<\beta<1),
\]
where $G$ is defined by \eqref{E:Gdef}.
Explicitly, for all $r$ with
$\phi^*(\alpha V_n r^p/V_q)>0$ and all $\alpha,\beta$ with
$0<\alpha<\beta<1$,
\[
\psi\Bigl(\frac{1}{\log(1/\beta)}\int_0^r
\log\frac{\phi^*(\alpha V_nx^p/V_q)}
{\phi^*(\alpha V_nr^p/V_q)}\,dx\Bigr)
\le
\frac{\phi^*((\alpha V_nr^p/V_q)-)}{\beta-\alpha}.
\]
Making the substitutions
\[
s^p:=\alpha V_nx^p/V_q,
\qquad
t^p:=\alpha V_nr^p/V_q,
\]
and using the fact that $\phi^*>0$ on $(0,m_p(X))$ (because
$\phi>0$ on $X$), we obtain
\[
\psi\Bigl(\frac{1}{\alpha^{1/p}\log(1/\beta)}
\Bigl(\frac{V_q}{V_n}\Bigr)^{1/p}
\int_0^t \log\frac{\phi^*(s^p)}{\phi^*(t^p)}\,ds\Bigr)
\le
\frac{\phi^*(t^p-)}{\beta-\alpha}
\]
for all $t\in(0,m_p(X)^{1/p})$ and all
$0<\alpha<\beta<1$. Since $n=p+q$, this is precisely
\eqref{E:Domarineq}. 
The definition of $\psi$ also immediately gives
\eqref{E:Ubound}.
\end{proof}

\end{document}
